%% kapa: exam=1 round=35 year=2024 period=1 form=A subject=civil_law subject_ko=민법 seq=6 number=6
%% source: Q-Net 감정평가사 기출문제 (cst00309) artlSeq=5213229 file=raw/1008/5213229/1교시 문제지 원본.pdf
%% conversion: kordoc 4.24.51 -> tools/gen1.py (auto); review=proofread-sonnet; flags=none
\begin{question}{6}
사단법인 A의 대표이사 甲이 A를 대표하여 乙과 금전소비대차계약을 체결하였다. 이에 관한 설명으로 옳지 않은 것은? (다툼이 있으면 판례에 따름)
\begin{choices}
  \item 甲이 A를 위하여 적법한 대표권 범위 내에서 계약을 체결한 경우, 그 계약의 효력은 A에게 미친다.
  \item 甲이 자신의 사익을 도모할 목적으로 대표권 범위 내에서 계약을 체결한 경우, 乙이 이 사실에 대해 알았다면 계약은 A에 대하여 효력이 없다.
  \item A의 정관에 甲이 금전소비대차계약을 체결할 수 없다는 규정이 있었지만 이를 등기하지 않은 경우, 乙이 이 사실에 대해 알았다면 A는 그 정관 규정으로 乙에게 대항할 수 있다.
  \item A의 乙에 대한 계약상 채무불이행책임 여부를 판단하는 경우, 원칙적으로 A의 고의ㆍ과실은 甲을 기준으로 결정한다.
  \item 만약 계약의 체결이 甲과 A의 이해가 상반하는 사항인 경우, 甲은 계약체결에 대해 대표권이 없다.
\end{choices}
\end{question}
\answer{3}
